\begin{table}[t]
\centering
\caption{Robustness Checks}
\label{tab:robust}
\begin{tabular}{lcccccc}
\hline\hline
 & \multicolumn{3}{c}{Quit Rate (per 100k)} & \multicolumn{3}{c}{Smoking Prevalence (\%)} \\
 & (1) & (2) & (3) & (4) & (5) & (6) \\
\hline
Grant$_z$ $\times$ Post & 173.5$^{**}$ & 412.8$^{***}$ & 281.2$^{***}$ & $-$0.597$^{***}$ & $-$0.114 & $-$0.044 \\
 & (82.4) & (111.4) & (87.4) & (0.118) & (0.112) & (0.204) \\
\hline
LA FE & Yes & Yes & Yes & Yes & Yes & Yes \\
Year FE & Yes & Yes & Yes & Yes & Yes & Yes \\
LA-specific trends & No & Yes & No & No & No & Yes \\
Baseline outcome $\times$ trend & No & No & No & No & Yes & No \\
Sample & Full & Full & 2013--19 & Full & Full & Full \\
\hline\hline
\end{tabular}
\begin{tablenotes}
\item \textit{Notes:} Each column reports a separate regression. Columns (1)--(3): CO-validated quit rate per 100,000. Column (1) is the baseline specification. Column (2) adds LA-specific linear trends. Column (3) excludes COVID years (2020--2022). Columns (4)--(6): smoking prevalence among adults 18+. Column (5) controls for baseline smoking prevalence interacted with a linear trend; column (6) adds LA-specific trends. The quit rate effect is robust: baseline grants are nearly uncorrelated with baseline quits ($r = 0.08$). The smoking prevalence effect reflects convergence ($r = 0.49$ between baseline grant and baseline smoking). Standard errors clustered at LA level in parentheses. $^{***}p<0.01$, $^{**}p<0.05$, $^{*}p<0.1$.
\end{tablenotes}
\end{table}
